\chapter{Microalbumin/Creatinine Ratio}

\section*{What Is This Test?}
The urine microalbumin/creatinine ratio, now properly called the urine albumin-to-creatinine ratio (UACR), measures small amounts of the protein albumin in the urine relative to the urine creatinine concentration. Unlike standard urine dipstick protein tests, which detect albumin only at concentrations greater than 300--500 mg/L, the UACR detects albumin at much lower concentrations (30--300 mg/g creatinine), previously termed ``microalbuminuria.'' The term ``microalbuminuria'' is being replaced by ``moderately increased albuminuria'' (A2) per KDIGO guidelines \cite{kdigo2012clinical}. Albumin is a 65 kDa protein that is normally almost entirely retained by the glomerular filtration barrier (composed of fenestrated endothelium, glomerular basement membrane with negatively charged heparan sulfate proteoglycans, and podocyte slit diaphragms). Small amounts of albumin (<30 mg/g creatinine) do pass the barrier but are almost completely reabsorbed by proximal tubular cells via megalin-cubilin receptors. The presence of 30--300 mg/g creatinine of albumin in the urine (A2) indicates early glomerular damage, most notably the earliest detectable stage of diabetic nephropathy. The test is performed on a random (spot) urine sample and is reported as milligrams of albumin per gram of creatinine (mg/g) or milligrams of albumin per millimole of creatinine (mg/mmol). The UACR is a cornerstone of CKD screening, staging, and cardiovascular risk stratification.

\section*{Alternative Names}
UACR; Urine Albumin-to-Creatinine Ratio; Microalbumin; Microalbuminuria Screen; Urine Microalbumin; ACR; Albumin/Creatinine Ratio; MA/Cr; Urine Albumin; Moderately Increased Albuminuria

\section*{Principle}
Urine albumin is measured by immunochemical methods, most commonly immunoturbidimetry or immunonephelometry \cite{tietz2012textbook}. Anti-human albumin antibodies react with albumin in the urine sample to form antigen-antibody complexes that scatter light. The amount of light scattered is proportional to the albumin concentration. The assay is highly sensitive (detection limit ~2--5 mg/L) and specific for albumin (does not cross-react with other urine proteins). Traditional urine dipsticks are insensitive below 150--300 mg/L and cannot detect ``microalbuminuria.'' Urine creatinine is measured simultaneously on the same sample by the enzymatic or Jaffe method. The albumin concentration (mg/L) is divided by the creatinine concentration (g/L or mmol/L) to yield the ratio. This correction for urine concentration eliminates the variability of a random spot collection compared to a timed (24-hour) collection \cite{wallach2022interpretation}. First-morning void urine samples are preferred because they are more concentrated and less affected by orthostatic proteinuria and exercise. This test is NOT performed by hematology analyzers.

\section*{History}
The association between proteinuria and kidney disease was recognized by Richard Bright in 1827 (Bright's disease). The concept of microalbuminuria was first described by Harry Keen and colleagues in 1963 using a radioimmunoassay. The landmark studies by Mogensen in the 1980s established microalbuminuria as the earliest sign of diabetic nephropathy and a predictor of progression to overt proteinuria and ESRD \cite{mogensen1984microalbuminuria}. The DCCT (Diabetes Control and Complications Trial, 1993) and UKPDS (United Kingdom Prospective Diabetes Study, 1998) demonstrated that intensive glycemic and blood pressure control reduced the development and progression of microalbuminuria \cite{ukpds1998intensive}. The KDIGO 2012 guidelines integrated albuminuria staging (A1, A2, A3) with eGFR staging into the CKD classification system (CGA classification: Cause, GFR, Albuminuria).

\section*{Why This Test Is Ordered}
\begin{itemize}
  \item Annual screening for diabetic nephropathy: ADA recommends annual UACR screening for patients with type 2 diabetes starting at diagnosis, and for patients with type 1 diabetes starting 5 years after diagnosis \cite{ada2023standards}
  \item Screening and risk stratification for chronic kidney disease: KDIGO recommends UACR for all patients at risk for CKD (diabetes, hypertension, cardiovascular disease, family history of CKD, age >60) \cite{levey2015glomerular}
  \item Cardiovascular risk assessment: albuminuria (even A2, 30--300 mg/g) is an independent predictor of cardiovascular events and mortality, even in patients with normal eGFR
  \item Monitoring response to therapy in diabetic nephropathy: ACE inhibitors and ARBs reduce albuminuria by 30--50\%, and the degree of albuminuria reduction correlates with renoprotection \cite{dezeeuw2004albuminuria}
  \item Detection of preeclampsia in pregnant women with hypertension: UACR >30 mg/mmol (or protein:creatinine ratio >30 mg/mmol) is suggestive of preeclampsia
  \item Screening for early kidney involvement in hypertension, SLE, and other systemic conditions
  \item Prognosis assessment before and after kidney transplantation
\end{itemize}

\section*{Primary vs Secondary Test}
The UACR is a primary screening test for diabetic nephropathy and a core component of CKD risk stratification per KDIGO guidelines. It is as essential as eGFR in CKD evaluation and should be checked annually in all diabetic patients and patients with hypertension or CKD risk factors.

\section*{How This Test Is Performed}
A random (spot) urine sample is collected, preferably the first-morning void, which correlates best with 24-hour albumin excretion and is less affected by orthostatic (postural) proteinuria and exercise. A midstream clean-catch urine is collected in a sterile container. No preservatives are needed. If a first-morning sample cannot be obtained, a random sample is acceptable, but results >30 mg/g should be confirmed with a first-morning void or a timed collection. The sample is refrigerated if not analyzed within 2 hours. Albumin is measured by immunoturbidimetry or immunonephelometry; creatinine is measured by enzymatic or Jaffe method. Results are reported as the ratio (mg albumin/g creatinine). Results are available within 1--2 hours.

\section*{What Preparation Is Needed}
No specific patient preparation required. However, several factors affect UACR and should be minimized: (1) Vigorous exercise within 24 hours can transiently increase albumin excretion (exercise-induced proteinuria). UACR should not be measured within 24 hours of strenuous exercise. (2) First-morning void is preferred to minimize orthostatic effects and maximize concentration. (3) Urinary tract infection causes transient albuminuria from inflammation---UACR should not be measured during active UTI; wait until infection resolves. (4) Fever, congestive heart failure, marked hypertension, and marked hyperglycemia (glucose >300 mg/dL) can transiently increase albuminuria independent of underlying nephropathy. (5) Menstruating women should delay testing until after menstruation to avoid blood contamination. (6) ACE inhibitors, ARBs, and SGLT2 inhibitors reduce albuminuria (therapeutic effect). NSAIDs may also affect results. (7) For diagnosis of persistent albuminuria, at least two out of three specimens over a 3--6 month period should be abnormal, because UACR has significant day-to-day variability (30--50\% coefficient of variation).

\section*{Contraindications and Precautions}
No absolute contraindications. Important considerations: (1) Day-to-day variability of UACR is high (30--50\%). A single elevated UACR should be confirmed with at least two additional measurements over 3--6 months before diagnosing persistent albuminuria. (2) Transient elevations: UTI, fever, heart failure, marked hyperglycemia, exercise, acute illness, menstrual blood contamination, orthostatic proteinuria. These should be excluded or resolved before diagnostic UACR measurement. (3) The test may be falsely elevated in very dilute urine with very low creatinine concentration, which can produce an artificially high ratio. If the urine creatinine is <30 mg/dL or >300 mg/dL, a timed (24-hour) collection may be more accurate. (4) False negatives: ACE inhibitors/ARBs reduce albuminuria (therapeutic effect). If UACR normalizes on these medications, it does not mean nephropathy has resolved---treatment is working. (5) UACR does not detect non-albumin proteins (e.g., Bence Jones proteins in multiple myeloma, tubular proteins in Fanconi syndrome). If non-albumin proteinuria is suspected, a urine protein-to-creatinine ratio (UPCR) or urine protein electrophoresis is required. (6) Pregnancy: UACR >30 mg/mmol is abnormal and may indicate preeclampsia. (7) The UACR is a screening and monitoring tool, not a substitute for 24-hour urine collection when precise quantification is needed (e.g., for nephrotic syndrome diagnosis: >3.5 g/day). (8) Commonly used cutoffs differ slightly by guideline: ADA uses <30 mg/g; KDIGO uses A1 <30 mg/g (normal to mildly increased), A2 30--300 mg/g (moderately increased), A3 >300 mg/g (severely increased).

\section*{Risks}
No risk to the patient. Urine collection is non-invasive. No blood draw is required.

\section*{What Happens After the Test}
Results available within 1--2 hours. KDIGO classification: A1 (<30 mg/g): normal to mildly increased; A2 (30--300 mg/g): moderately increased (formerly microalbuminuria); A3 (>300 mg/g): severely increased (formerly macroalbuminuria/overt proteinuria). The combination of eGFR and albuminuria categories (the KDIGO heat map) provides risk stratification for CKD progression, ESRD, cardiovascular events, and mortality: Patients with UACR >300 mg/g and eGFR <30 have the highest risk. For diabetic nephropathy: A2 (30--300 mg/g) indicates incipient nephropathy---this is the stage where intervention is most effective. Without intervention, approximately 40--50\% of type 1 diabetics with A2 progress to overt nephropathy (A3) over 10--15 years; with aggressive blood pressure control (target <130/80), ACE inhibitor/ARB therapy, glycemic control (HbA1c <7\%), and SGLT2 inhibitor use, progression can be slowed or halted. A 50\% reduction in UACR with therapy correlates with a 50\% reduction in ESRD risk. A value >300 mg/g indicates overt diabetic nephropathy with significant risk of progression to ESRD.

\section*{Understanding Results}

\subsection*{Below Normal}
\begin{itemize}
  \item \textbf{UACR <30 mg/g:} Normal to mildly increased (A1). Indicates intact glomerular filtration barrier. Does not exclude renovascular disease, hypertensive nephrosclerosis (which may not cause significant albuminuria early), or chronic interstitial nephritis.
  \item \textbf{Normalization in response to therapy:} ACE inhibitors/ARBs/SGLT2 inhibitors reduce albuminuria. UACR dropping from A2 to A1 on therapy indicates treatment response and reduced risk of progression.
  \item \textbf{Non-albumin proteinuria with normal UACR:} Multiple myeloma (Bence Jones proteins---light chains), some forms of tubulointerstitial disease. Urine protein electrophoresis or UPCR would be abnormal despite normal UACR.
\end{itemize}

\subsection*{Above Normal}
\begin{itemize}
  \item \textbf{Moderately increased albuminuria (A2): UACR 30--300 mg/g.} Earliest clinically detectable stage of glomerular injury.
  \item \textbf{Diabetic nephropathy (most common cause worldwide):} Hyperglycemia causes glomerular hyperfiltration, basement membrane thickening, mesangial expansion, and podocyte injury. A2 typically develops after 5--10 years of type 1 diabetes; may be present at diagnosis in type 2 diabetes (which often goes undetected for years). With intensive glycemic and blood pressure control, A2 may regress to A1 in 30--50\% of patients.
  \item \textbf{Hypertensive nephrosclerosis:} Chronic hypertension causes glomerular capillary hypertension and hyperfiltration injury, leading to albuminuria and progressive nephron loss. A2 develops in 15--30\% of hypertensive patients, particularly African Americans. ACE inhibitors or ARBs are first-line therapy.
  \item \textbf{Early glomerulonephritis:} IgA nephropathy, membranous nephropathy, FSGS, minimal change disease (though MCD typically presents with nephrotic-range proteinuria, early disease may show A2). Lupus nephritis (SLE) requires regular UACR screening---A2 in SLE warrants kidney biopsy.
  \item \textbf{Cardiovascular disease:} A2 is an independent risk factor for myocardial infarction, stroke, and cardiovascular mortality, even with normal eGFR. This is thought to reflect systemic endothelial dysfunction, of which albuminuria is a renal manifestation.
  \item \textbf{Severely increased albuminuria (A3): UACR >300 mg/g.} Indicates significant glomerular damage with overt proteinuria. Corresponds to 24-hour urine protein >300--500 mg/day. Risk of progression to ESRD is substantially elevated. In diabetes, A3 indicates overt diabetic nephropathy---approximately 50\% of patients with A3 progress to ESRD within 10 years without aggressive intervention. Nephrotic syndrome is defined as >3.5 g/day total protein, which is >3,500 mg/g UACR (approximate).
\end{itemize}

\section*{Normal Results}
\begin{itemize}
  \item \textbf{A1 (normal to mildly increased):} <30 mg/g creatinine (<3.4 mg/mmol)
  \item \textbf{A2 (moderately increased, formerly microalbuminuria):} 30--300 mg/g (3.4--34 mg/mmol)
  \item \textbf{A3 (severely increased, formerly macroalbuminuria):} >300 mg/g (>34 mg/mmol)
  \item \textbf{First-morning void preferred:} Random sample reference ranges are the same but confirmation required
  \item \textbf{Timed 24-hour urine albumin:} A2 = 30--300 mg/24h; A3 = >300 mg/24h
  \item \textbf{Nephrotic range proteinuria:} >3,500 mg/24h total protein (~3,500 mg/g on UPCR)
  \item \textbf{KDIGO CGA classification:} Combined eGFR + Albuminuria category (A1/A2/A3)
\end{itemize}

\section*{Follow-up Tests}
\begin{itemize}
  \item \textbf{Repeat UACR (first-morning void):} Confirm abnormal result. At least 2 of 3 specimens over 3--6 months must be abnormal to diagnose persistent albuminuria.
  \item \textbf{Urine protein-to-creatinine ratio (UPCR):} Measures total urine protein (not just albumin). Useful when A3 is present for quantification, or when non-albumin proteinuria is suspected (multiple myeloma).
  \item \textbf{24-hour urine collection for total protein and creatinine clearance:} Gold standard quantitation. Used when UACR is inconsistent or precise measurement is needed (nephrotic syndrome diagnosis, pregnancy).
  \item \textbf{Serum creatinine and eGFR:} Essential for complete CKD staging. Combined eGFR + UACR provides the KDIGO risk category.
  \item \textbf{Renal ultrasound:} Assess kidney size, cortical thickness, echogenicity, and rule out obstruction.
  \item \textbf{Serum albumin and lipid panel:} In nephrotic-range proteinuria, albumin falls and lipids rise (nephrotic syndrome).
  \item \textbf{HbA1c and blood glucose:} Glycemic control assessment in diabetic nephropathy.
  \item \textbf{Urine microscopy:} Assess for dysmorphic RBCs, RBC casts (glomerulonephritis), WBC casts (interstitial nephritis).
  \item \textbf{Serologic workup:} ANA, anti-dsDNA, ANCA, anti-GBM, complement C3/C4, hepatitis B/C, HIV for glomerulonephritis evaluation.
  \item \textbf{Kidney biopsy:} Indicated if UACR >300 mg/g with unclear etiology, rapidly progressive, or with concomitant hematuria and RBC casts.
\end{itemize}

\section*{References}
\bibliographystyle{unsrtnat}
\bibliography{references}

\clearpage
